\section{Turing Machines (Máquinas de Turing)}
\label{sec:modeling-turing}

A categoria \textbf{Turing Machines} permite modelar processos computacionais através de máquinas abstratas que manipulam símbolos numa fita.

\subsection{Objetos}
Os objetos desta categoria são máquinas de Turing (determinísticas) finitas. Uma máquina de Turing $M = (Q, \Gamma, \delta, q_0, F, B)$ é codificada em MATLAB/Octave como uma \textit{struct} composta pelos seguintes elementos:
\begin{itemize}
    \item \texttt{Q} --- número de estados ($|Q|$), identificados por inteiros $1, \ldots, |Q|$;
    \item \texttt{Gamma} --- vetor de símbolos de fita (o símbolo em branco é o índice 1 por convenção);
    \item \texttt{delta} --- matriz de transição de dimensão $|Q| \times |\Gamma|$, onde \texttt{delta(q, a)} é um vetor $[q', a', d]$ com: $q'$ = novo estado, $a'$ = símbolo escrito, $d \in \{-1, 0, 1\}$ = direção;
    \item \texttt{q0} --- estado inicial;
    \item \texttt{F} --- vetor de estados finais/aceitação.
\end{itemize}

\begin{lstlisting}[language=Matlab, caption={Representação de uma Máquina de Turing}]
% Exemplo: MT que aceita a^n b^n (simplificado)
M.Q     = 5;              % estados 1..5
M.Gamma = [0, 1, 2, 3];  % 0=branco, 1='a', 2='b', 3='X'
M.blank = 1;              % indice do branco em Gamma
M.q0    = 1;              % estado inicial
M.F     = [5];            % estados de aceitacao

% delta(q, sym) = [new_q, new_sym, dir]   dir: -1=L, 0=N, 1=R
M.delta = zeros(M.Q, numel(M.Gamma), 3);
% Transicoes (exemplo parcial):
M.delta(1, 2, :) = [2, 4, 1];  % q1, 'a' -> escreve 'X', vai direita, q2
M.delta(2, 2, :) = [2, 2, 1];  % q2, 'a' -> fica 'a', vai direita, q2
M.delta(2, 3, :) = [3, 4, -1]; % q2, 'b' -> escreve 'X', vai esquerda, q3
% ... (restantes transicoes omitidas por brevidade)
\end{lstlisting}

\subsection{Morfismos}
Um morfismo $f: M_1 \to M_2$ é um par de mapas $(f_Q, f_\Gamma)$ --- sobre os conjuntos de estados e de símbolos de fita, respetivamente --- que preservam a função de transição, o estado inicial, os estados finais, e o símbolo em branco. Formalmente, devem ser satisfeitas as seguintes condições:
\begin{enumerate}
    \item $f_Q(q_0^{(1)}) = q_0^{(2)}$ (preservação do estado inicial);
    \item $f_Q(F_1) \subseteq F_2$ (preservação dos estados finais);
    \item $f_\Gamma(B_1) = B_2$ (preservação do símbolo em branco);
    \item Para cada $(q, a) \in Q_1 \times \Gamma_1$, se $\delta_1(q, a) = (q', a', d)$, então $\delta_2(f_Q(q), f_\Gamma(a)) = (f_Q(q'), f_\Gamma(a'), d)$.
\end{enumerate}

\begin{lstlisting}[language=Matlab, caption={Verificação de um morfismo de Máquinas de Turing}]
function ok = is_TM_morphism(M1, M2, fQ, fGamma)
    % fQ: vetor de tamanho M1.Q  (fQ(q) em {1..M2.Q})
    % fGamma: vetor de tamanho numel(M1.Gamma)
    ok = true;
    
    % 1) Estado inicial
    if fQ(M1.q0) ~= M2.q0
        ok = false; return;
    end
    
    % 2) Estados finais
    for q = M1.F
        if ~ismember(fQ(q), M2.F)
            ok = false; return;
        end
    end
    
    % 3) Simbolo em branco
    if fGamma(M1.blank) ~= M2.blank
        ok = false; return;
    end
    
    % 4) Funcao de transicao
    for q = 1:M1.Q
        for a = 1:numel(M1.Gamma)
            tr = squeeze(M1.delta(q, a, :))';
            q2 = tr(1); a2 = tr(2); d = tr(3);
            tr2 = squeeze(M2.delta(fQ(q), fGamma(a), :))';
            if tr2(1) ~= fQ(q2) || tr2(2) ~= fGamma(a2) || tr2(3) ~= d
                ok = false; return;
            end
        end
    end
end
\end{lstlisting}